../cictikz/src/cictikz/data/tex/cictikz_preamble.tex

% Third step of the R-2R build-up: hang a 2R leg on V_1 as well, and the
% ladder repeats — 2R in parallel with the 2R that dac_2r_1 reduced to is R
% again, which is why every further section is the same drawing.
%
% Each leg carries half the current of the one to its left: I_1 = V_1/2R and
% I_0 = V_0/2R with V_0 = V_1/2, which is the binary weighting the DAC needs.

\begin{circuitikz}[american, thick, transform shape, circuit ee IEC]

  ../cictikz/src/cictikz/data/tex/cictikz_lib.tex
  ../cictikz/src/cictikz/data/tex/rdac_lib.tex

  \def\yr{3.0}

  \coordinate (v1) at (0,\yr);
  \draw (v1) \hresistor{$R$} coordinate (v0);
  \node[anchor=south east] at ($(v1)+(0.1,0.15)$) {$V_1$};
  \node[anchor=south] at ($(v0)+(0,0.15)$) {$V_0$};

  % ---- 2R across to the summing node ----
  \draw (v0) \hresistor{$2R$} -- ++(0.7,0) coordinate (vg);
  \rdacVgnd{(vg)}

  % ---- A 2R leg on each node ----
  \draw ($(v1)+(0,-1.6)$) \vresistor{$2R$};
  \draw ($(v1)+(0,-1.6)$) -- ++(0,-0.3) \vground;
  \draw ($(v0)+(0,-1.6)$) \vresistor{$2R$};
  \draw ($(v0)+(0,-1.6)$) -- ++(0,-0.3) \vground;

  % ---- The leg currents, as the original marks them ----
  \draw[-{Latex[length=3mm]}] ($(v1)+(-0.7,-0.5)$) -- ++(0,-1.0)
    node[anchor=east, midway, xshift=-2mm] {$I_1$};
  \draw[-{Latex[length=3mm]}] ($(v0)+(1.3,-0.5)$) -- ++(0,-1.0)
    node[anchor=west, midway, xshift=2mm] {$I_0$};

\end{circuitikz}
\end{document}
